\begin{textsample}{POS Dim 4 – persona\_gemini – Score 16.00 – t911\_gemini.txt}  \label{ex:f4_pos_045}
As melting Arctic \textbf{ice}, a direct result of climate change, opens up the \textbf{region}, oil companies and \textbf{fishing} \textbf{fleets} are able to move in and exploit it.
The \textbf{ecosystem} around Svalbard is particularly at risk from destructive \textbf{fishing} practices. \textbf{Bottom} trawlers are a major \textbf{threat}, destroying \textbf{species} that live on the \textbf{seabed} like \textbf{sea} pens and corals.
This \textbf{activity} runs the risk of eradicating \textbf{species} that have not yet even been discovered.
At the same time, bycatch from \textbf{fishing} \textbf{operations} harms near-threatened \textbf{species} such as the Greenland \textbf{shark}.
Discarded plastic \textbf{fishing} gear creates another hazard, entangling polar bears.
The impact extends beyond physical harm.
Underwater noise from \textbf{vessels} disrupts narwhals and beluga \textbf{whales}, affecting their ability to communicate and breed.
The \textbf{fishing} \textbf{fleets} also reduce the amount of food available for walruses.
This adds to the stress these \textbf{animals} already experience from the loss of \textbf{sea} \textbf{ice}, demonstrating how these pressures compound in an interconnected \textbf{ecosystem}.
We must urge the Norwegian government to protect Svalbard from this destructive \textbf{fishing}.

% matched lemmas: activity, animal, bottom, ecosystem, fishing, fleet, ice, operation, region, sea, seabed, shark, specie, threat, vessel, whale
\end{textsample}
